\scanpage{17}
$A\times_S U\simeq A_U$ can (\unsure{$=$|même|$\Leftarrow$} raison). De plus, je dis qu'on a un isom $T_p(A)\simeq M$, qui recolle les $T_p(A_{\deleted{\unsure{s}}(s)})\simeq M_{(s)}$. En effet, pr tt entier $n$, par le résultat déjà signalé de pl. fid., on aura un isom $\deleted{T_p(A}{}_{p^n}A\simeq M(n)$ qui recolle les
\diagram{\begin{tikzcd}[ampersand replacement=\&,column sep=small,row sep=small]
{}_{p^n}(A_s) \arrow[d,equal] \arrow[r,"\sim"] \& \deleted{\unsure{p}}M_{(s)}(n) \arrow[d,equal] \\
({}_{p^n}A)_s \& (M(n))_s
\end{tikzcd}}
. Sur $U$, c'est l'isom. tautologique. Ceci assure l'unicité.

Il reste à prouver la pleine fidélité précédente (on la laisse au lecteur), et le

\margin{Douteux, sans doute faux}
\uline{Lemme}: Un ``syst local de schémas abéliens'' sur un schéma [loc noeth] \deleted{\unsure{est}} provient d'un schéma abélien.
\margin{\unsure{et} [\hspace{.6em}] inutile?}

On peut supposer $S$ noethérien. \inserted{On est ramené à prouver que} \deleted{Soit} si $U$ \deleted{\unsure{le plus grand}} soit un ouvert tel que l'on ait un schéma abélien $A_U$ sur $U$, et $t$ un pt maximal de $S-U=T$, on peut prolonger sur un voisinage $V$ de $t$. Or \deleted{soit $B$ un} $V$ un voisinage ouvert de $t$, tel que $A_{(t)}$ provient d'un schéma abélien sur $V$, \deleted{et \deleted{soit} tel que $V$ ne contient pas les $x\in\operatorname{Ass}\mathcal O_S$ tels que $\overline{\{x\}}\not\ni t$,} \inserted{et tel que $B\mid U\cap V\simeq A_U$} \deleted{Je dis que pour tt $s\in V$, il \unsure{faut}} (isom compatible avec les isoms $B\times\operatorname{Spec}\mathcal O_{S,s}\simeq A_{(s)}$ déjà connus pour $s\in U\cap\operatorname{Spec}\mathcal O_{S,t}=U_{(s)}$). \deleted{Soit} On peut supposer que $x\in\operatorname{Ass}\mathcal O_S$, $\bar x\not\ni\{t\}\Rightarrow x\notin V$, et que $V\cap T\subset\overline{\{t\}}$ \deleted{\unsure{de}}. Soit alors $s\in V\cap T$, donc $s\in\overline{\{t\}}$
% archive sheet number: 16

